\section{Three-body clustering and cohomological response}
\label{sec:cohomology}

Figure~\ref{fig:mechanisms} compares three protected constructions for which
the response is nonzero: a three-body clustering kernel, a random
supercharge, and a local lattice supercharge.  The common feature is an
exactly degenerate zero-energy fiber separated from its complement.  The
origin of that fiber and the sampling measure are different.  We therefore
compare the algebra and the statistical result within each model before
placing their evidence cells next to one another.  In particular, a tangent
panel at one fixed Hamiltonian is not interchangeable with an independent
disorder realization.

\subsection{Moore--Read: an exact moving kernel}

The Moore--Read parent is a sum of positive three-particle constraints,
Eq.~\eqref{eq:moore-read-parent}.  The states in its kernel obey the same
three-body clustering condition at exactly zero energy
\cite{mooreread1991,zhang2023parent}.  For the fixed-four-quasihole sequence
$N_\phi=N+2$, the two opened systems have $(N,D)=(4,
\MooreReadNFourRank)$ and $(6,\MooreReadNSixRank)$.  The first positive
energies, \MooreReadNFourExternalGap\ and
\MooreReadNSixExternalGap, respectively, verify that the response is taken
for an isolated projector rather than for a numerically selected portion of
a continuum.

The guiding-density directions are an exactness-preserving generator
transport.  If $G_a$ is lifted to the $N$-particle Fock space by
$\Gamma_N(G_a)$, differentiation of the transported constraint gives
\begin{equation}
 (\partial_aH)P=-H\Gamma_N(G_a)P,
 \qquad
 X_a=Q\Gamma_N(G_a)P.
 \label{eq:mr-generator-response}
\end{equation}
Thus $S_a=P(\partial_aH)P=0$, while the off-fiber block $X_a$ is generally
nonzero.  This identity is checked against the parent action and complement
leakage for every stored panel.  It matters conceptually: the response is not
the derivative of a multiplet that immediately splits.  It is the derivative
of an exactly degenerate kernel that moves under a finite invertible
constraint path.

The finite-size calculations used two related, but inequivalent, fourth-order
statistics.  The historical two-pairing residual in
Fig.~\ref{fig:mechanisms}(a) subtracts a separable proper-complex Gaussian
prediction.  Its local-panel medians are
\MooreReadNFourLocalTwoPairRFour\ and
\MooreReadNSixLocalTwoPairRFour; the structured-line controls are
\MooreReadNFourStructuredTwoPairRFour\ and
\MooreReadNSixStructuredTwoPairRFour.  Those values describe the old
$R_4^{2\mathrm{pair}}$ observable.  They neither include the
pseudocovariance pairing nor reproduce an arbitrary entrywise covariance, so
they are not a test of the complete Gaussian law.

The v12 raw, uncentered three-pairing calculation is also retained only for
provenance.  Its response arrays were neither population-centered nor
channel-whitened before the moment subtraction.  The corresponding raw
estimates, \MooreReadNFourRawMomentResidualEstimate\ and
\MooreReadNSixRawMomentResidualEstimate, cannot be interpreted as connected
cumulants.  A nonzero raw value may contain the finite-population mean and is
therefore claim-ineligible even when a nominal normal-reference probability
is small.  Figure~\ref{fig:mechanisms}(b) marks this distinction explicitly:
``exact fiber and response'' and ``centered complete covariance'' are
separate evidence questions.

For the corrected $R_4^{\mathrm{full}}$ calculation, one fixed-base tangent
panel is the statistical unit.  Each system has a registered finite
population of torus guiding-density anchors, and linearity gives its exact
population mean before any panels are assigned to inference.  The mean
vanishes for the uniform traceless guiding-density population; it is not
estimated from the observed inference half.  The \FullRFourTrainingCount\
panels in the training half determine only the $2\times2$ channel Gram
matrix.  Its inverse square root is then frozen and applied to the disjoint
\FullRFourInferenceCount\ panels.  We refer to this as the frozen 12/12
split.  The training half alone fixes the whitening transformation.  All
three Wick pairings, including the pseudocovariance term, are subtracted from
the inference-panel U-statistic.

Uncertainty is obtained by deleting one complete inference panel and
recomputing the contraction.  The registered scalar direction is tested with
a Student-$t$ reference with \FullRFourStudentTDegreesOfFreedom\ degrees of
freedom.  A Bonferroni correction controls the family level
\FullRFourFamilyAlpha\ across five registered primary cases
(\FullRFourPrimaryFamilySize\ in total), giving per-case level
\FullRFourPerCaseAlpha.  This protocol
is conditional on exchangeability of the fixed-base panels; it is not an
estimate over Moore--Read Hamiltonians or over parent-Hamiltonian families.

At $N=4$, $D=\MooreReadNFourRank$, the directional estimate is
\MooreReadNFourFullRFourDirectionalEstimate, while the familywise one-sided
lower bound is \MooreReadNFourFullRFourFamilywiseLower.  The registered
positive-direction gate does not pass.  At $N=6$,
$D=\MooreReadNSixRank$, the corresponding estimate is
\MooreReadNSixFullRFourDirectionalEstimate\ and the familywise lower bound is
\MooreReadNSixFullRFourFamilywiseLower; this case passes.  The positive cell
in Fig.~\ref{fig:mechanisms}(c) is consequently a finite-rank $N=6$ result.
It is not an asymptotic statement, and the absence of a pass at $N=4$ forbids
a monotone scaling claim from these two ranks alone.  The reverse split is
descriptive and is not pooled with the primary split, even where its point
estimate has the same sign.

This result establishes a connected fourth-order response contraction at
one rank for the stated tangent population after complete covariance and
pseudocovariance subtraction.  It does not establish that all fourth-order
components are nonzero, that the effect persists with $D$, or that different
three-body parents share the same law.  It also does not identify a dynamical
Lyapunov exponent: the measured object is the geometry of an exactly
degenerate projector over coupling space.

There are two further reasons not to read the $N=4,6$ comparison as a fit.
First, the ambient occupation-space dimensions and the protected ranks change
together, so a change in the scalar contraction cannot be assigned uniquely
to $D$.  Second, the response channels are generated by a discrete torus
population whose geometry also changes with $N_\phi$.  An asymptotic test
would have to hold the operator prescription fixed in physical units, open
additional ranks, and preregister the scaling quantity before examining the
new outcomes.  Here the two cases instead test whether the same finite-panel
protocol can be carried out at two nontrivial ranks and whether its familywise
gate is resolved in either one.

The fixed-base qualification does not contradict exact transport.  Every
$G_a$ defines a finite isospectral constraint path, but the collection of
panels is evaluated at the same base parent and base projector.  The
jackknife therefore samples the registered distribution of tangent
directions through one point of the projector bundle.  It does not sample
different points along a path, and it does not absorb the correlation among
channels within one panel into the sample count.  This distinction is why
the full response tensor of a panel, rather than each generator or matrix
entry, is deleted in the uncertainty calculation.

\subsection{Random-supercharge cohomology}

The random supersymmetric model replaces a clustering ideal by a chain
complex.  A cubic nilpotent supercharge $Q$ maps a charge-$r$ sector to the
charge-$(r-3)$ sector, and
\begin{equation}
 H_r=Q_{r+3}Q_{r+3}^{\dagger}+Q_r^{\dagger}Q_r
 \label{eq:syk-hodge-laplacian}
\end{equation}
is its Hodge Laplacian.  Its zero modes are harmonic representatives of
$Q$-cohomology, $\ker Q_r/\operatorname{im}Q_{r+3}$, and hence are annihilated
by both $Q$ and $Q^\dagger$ \cite{fu2017susy}.  Exact degeneracy follows from
nilpotency and positivity, not from spatial locality.  This is the finite
supersymmetric sector in which non-Abelian Berry curvature was previously
proposed as a chaos probe \cite{chen2026}.

Let $\Pi_-=Q_{r+3}Q_{r+3}^{\dagger}H_r^+$ and
$\Pi_+=Q_r^{\dagger}Q_rH_r^+$ be the orthogonal projectors onto the exact and
coexact positive-energy subspaces, where $H_r^+$ is the Moore--Penrose inverse
of $H_r$ on the complement of its kernel.  Differentiating the two Hodge terms
then separates the response into $X_a^-=\Pi_-X_a$ and
$X_a^+=\Pi_+X_a$,
\begin{equation}
 X_a=X_a^{-}\oplus X_a^{+},
 \qquad
 \langle X_a^{-},X_b^{+}\rangle_{\mathrm{HS}}=0.
 \label{eq:hodge-response-split}
\end{equation}
Here the direct resolvent response and the sum of the two branch responses
are checked independently.  The branch weights, channel covariances, and
left/right covariance spectra can therefore be used to build a
Hodge-resolved Gaussian comparison without opening the fourth-order outcome.
This branch structure is physical information imposed by the complex; it
should not be discarded by replacing the full response with an isotropic
rectangular matrix.

The registered $N=14$ test used a complete independent disorder realization
as the uncertainty unit.  The adjacent- and central-charge sparse panels have
observed median historical residuals
\SusySykNFourteenMedianOne\ and
\SusySykNFourteenMedianTwo.  Before those outcomes were unsealed, the null
distribution was generated from the stored two-point Hodge signatures.  The
held-out medians lie outside both the collapsed and Hodge-resolved prediction
intervals.  This rejects the frozen separable proper-complex Gaussian null
for the registered response statistic.

The qualifier ``separable'' is essential.  That null matches marginal
channel, target, external, and exact/coexact covariance spectra, but it is not
the complete entrywise covariance $C$ together with pseudocovariance $\Pi$.
Consequently, its rejection can arise from a connected fourth cumulant or
from nonseparable two-point structure omitted by the null.  The random-SYK
cell for centered $R_4^{\mathrm{full}}$ remains not tested.  We do not use the
historical residual to claim complete-covariance non-Gaussianity, and we do
not pool its independent-realization bootstrap with the fixed-base panel
jackknife used for Moore--Read.

This separation also identifies the useful next test.  Random couplings
supply genuine Hamiltonian-to-Hamiltonian sampling, whereas the Moore--Read
calculation supplies an exact moving-kernel path at one parent.  A common
claim would require the same centered full-covariance observable and a
matched tangent prescription in both ensembles.  Those cells are absent from
the present registry.

The Hodge split also changes the meaning of apparent isotropy.  Unitary
rotations within the exact image and within the coexact image may leave all
stored covariance spectra invariant while changing correlations between the
two embedded subspaces and the protected fiber.  A collapsed null that keeps
only the total left and right spectra forgets even more information.  The
sealed v7 comparison was designed to ask whether these progressively richer
separable nulls predicted the held-out median.  Their failure is reproducible
evidence for covariance memory beyond those summaries; it is deliberately
not relabeled as the complete cumulant of Sec.~\ref{sec:statistics}.

\subsection{Local lattice cohomology}

The lattice construction retains the Hodge mechanism but changes both
locality and sampling.  On a union of $m$ six-site cycles, the nilpotent
supercharge acts on the independence complex, and its harmonic sector has
ranks $D=\LatticeSusyMOneRank$, \LatticeSusyMTwoRank, and
\LatticeSusyMThreeRank\ for $m=1,2,3$.  The corresponding positive gaps are
\LatticeSusyMOneExternalGap, \LatticeSusyMTwoExternalGap, and
\LatticeSusyMThreeExternalGap.  Local changes of the site couplings move the
harmonic projector.  The response again splits into exact and coexact
branches; their measured local-panel balance factors are
\LatticeSusyMOneHodgeBalance, \LatticeSusyMTwoHodgeBalance, and
\LatticeSusyMThreeHodgeBalance.  Thus the local family is neither a
fixed-projector control nor a single-branch rectangular ensemble.

Panel (a) shows the historical two-pairing residual for local and isotropic
tangents.  The local values
\LatticeSusyMOneLocalTwoPairRFour,
\LatticeSusyMTwoLocalTwoPairRFour, and
\LatticeSusyMThreeLocalTwoPairRFour\ differ from their isotropic counterparts
\LatticeSusyMOneIsotropicTwoPairRFour,
\LatticeSusyMTwoIsotropicTwoPairRFour, and
\LatticeSusyMThreeIsotropicTwoPairRFour.  These finite-size differences show
that the accessible tangent class matters.  As in Moore--Read, however, the
old $R_4^{2\mathrm{pair}}$ numbers do not answer the complete-covariance
question.  The v12 raw estimates are likewise uncentered and are not used for
a claim.

The corrected calculation uses the exact population mean obtained by
enumerating the registered ordered local-coordinate marginals.  This mean is
nontrivial because the discrete local-tangent construction is not symmetric
entry by entry.  After population centering, the same frozen 12/12 split,
training-only channel whitening, three-pairing U-statistic, Student-$t$
reference, and five-case Bonferroni family are used.  This alignment makes the
Moore--Read and lattice rows comparable at the level of the estimator, but it
does not turn the lattice panels into independent graph realizations.

Disjoint cycle components make the algebraic rank sequence exactly auditable,
but they also impose a tensor-product ancestry on the cohomology.  Increasing
$m$ simultaneously adds degrees of freedom and repeats a local building
block.  The sequence is therefore useful for checking how a known branch
structure enters the response and the estimator; it is not a connected-
lattice thermodynamic sequence.  A claim of locality-driven universality
would require connected graphs, a matched local tangent density, and an
independent ensemble at each size.  None of those missing conditions can be
replaced by treating the many coordinate pairs at one graph as independent
Hamiltonians.

The primary directional estimates are
\LatticeSusyMOneFullRFourDirectionalEstimate,
\LatticeSusyMTwoFullRFourDirectionalEstimate, and
\LatticeSusyMThreeFullRFourDirectionalEstimate.  Their familywise lower
bounds are \LatticeSusyMOneFullRFourFamilywiseLower,
\LatticeSusyMTwoFullRFourFamilywiseLower, and
\LatticeSusyMThreeFullRFourFamilywiseLower.  None of the three registered
positive-direction gates passes; failure of the positive-direction gate does
not prove Gaussianity.  It supplies neither an equivalence test nor a
two-sided upper bound on every component of $\mathcal K$.  A negative point
estimate in a preselected direction may coexist with a nonzero cumulant in a
different direction, and a small finite panel family may lack power.  The
reverse split is descriptive here as well.

The appropriate conclusion is therefore asymmetric.  Local lattice
cohomology supports an exact, moving, branch-reducible projector with
nontrivial response, but the registered positive full-cumulant contraction is
not resolved for $D=\LatticeSusyMOneRank,
\LatticeSusyMTwoRank,\LatticeSusyMThreeRank$.  Moore--Read resolves that
contraction only in the displayed rank-\MooreReadNSixRank\ case.  Random
supercharge SYK rejects a narrower separable null using independent disorder
realizations, while its complete entrywise-covariance cell remains open.
These are three different statements, not three measurements of a universal
number.
